% Pore width from porosity and thread width
\[
\varepsilon = \frac{s^2}{(s + w)^2} \quad \Rightarrow \quad s = \frac{w \sqrt{\varepsilon}}{1 - \sqrt{\varepsilon}}
\]

% Single pore area
\[
A_{\text{pore}} = s^2
\]

% Number of pores
\[
N = \frac{A_{\text{pores,total}}}{A_{\text{pore}}} = \frac{\varepsilon \, A_{\text{total}}}{s^2}
\]

Sure! The **solid mass** in a porous mesh can be calculated from total area, thickness, porosity, and material density. The LaTeX code is:

```latex id="g1k7zo"
% Solid mass in a porous material
\[
m_{\text{solid}} = \rho_{\text{material}} \, A_{\text{total}} \, t \, (1 - \varepsilon)
\]

% Where:
% \rho_{\text{material}} = material density (kg/m^3)
% A_{\text{total}} = total area of mesh (m^2)
% t = thickness of the mesh (m)
% \varepsilon = porosity (fraction)
```


